\section{Threats to Validity}\label{threats-to-validity}

\textbf{Construct validity.} Terminal names are not original Jira status-category metadata, and no human workflow ground truth or inter-rater agreement is available. A recorded terminal state may include administrative closure rather than delivered functionality. Archived Start\_Date and End\_Date have no historical revision log in this dataset; their availability and stability at sprint start are assumptions. Project ownership is likewise snapshot metadata. These limitations prevent an unconditional claim of historically faithful business-commitment prediction.

\textbf{Selection and missingness.} Only 77.04\% of reconstructed commitments receive an assessable outcome, and eligibility further restricts evaluation to 14 projects. Unassessable outcomes may systematically differ from retained outcomes. Retrospective inconsistency filtering also uses later quality evidence. Static attributes missing timestamped support are not backfilled, potentially weakening the start-time comparator relative to complete real-world initial issue data. The experiment measures the value of the reconstructed information sets, not all information a team originally possessed.

\textbf{Internal validity.} We audit event prefixes, freeze partitions, purge unavailable endpoint labels, and keep preprocessing and recalibration outside test fitting. These measures do not establish complete changelog capture or detect every incorrectly timestamped record. Adding dynamic variables changes representation dimension as well as observations; the zero-percent control helps expose, but cannot eliminate, fitting effects. One fixed CatBoost configuration and one seed do not establish algorithmic optimality.

\textbf{Statistical conclusion validity.} Sprint pairing preserves the primary comparison but does not eliminate dependence from recurring issues, overlapping work, or repositories shared across projects. Cross-project bootstrap has only 14 project units and overlapping training populations across folds. Supplementary comparisons are not multiplicity-adjusted and are not confirmatory replacements for H1. Calibration plots have bins of unequal sample size; small-bin apparent deviations need cautious interpretation.

\textbf{External validity.} The sample concerns historical open-source Jira data rather than current enterprise deployment. Project holdout is not organization holdout, and cross-project evaluation is not a global time-cutoff experiment. Evaluated alert cohorts contain only assessable outcomes, whereas an online system would not know outcome assessability in advance. Offline recall, precision, and lead time do not demonstrate causal benefit from alert intervention.
